\section{Vertex filtrations and uniqueness}\label{sec:matchinguniqueness}

We first construct the local residue data at the vertices. We then prove
a converse: a finite regular harmonic path satisfying first-order matching
is the variational minimum. Both arguments use the full norm motion,
rather than only its unordered velocity spectrum.

\subsection{The residue representation at a vertex}

\begin{proposition}\label{prop:vertexfiltrations}
Let $y$ be a finite regular harmonic path satisfying
\eqref{eq:matching}. At every support vertex it determines an
actual reduced $A_2$ representation with two increasing-phase
filtrations. Their quotients are the incident incoming and
outgoing edge objects, with their actual arrows. These
identifications extend along edges and remove artificial
subdivisions. Changes of apartment bases give isomorphic data.
\end{proposition}

\begin{proof}
Fix a subdivision time $u_0$, put $s=u-u_0$, and choose
$y_v(u_0)$-orthonormal reference bases at the two vertices.
Normalize the left and right apartment frames in the same way.
Every transition matrix $P$ lies in $\GL_n(\OO)$, with
invertible reduction $\overline P\in\GL_n(\CC)$.
The matrix $\overline P^{-1}P$ is $I+N$, where all nonzero
entries of $N$ have positive valuation. For the finitely many
matrices and inverses choose a common positive lower bound
$\delta$ for those valuations. As in
Lemma~\ref{lem:taildescent}, replacing $P$ by $\overline P$
preserves exactly every weighted norm whose weight spread is
less than $\delta$. Our one-sided weights tend to zero.
Thus, on a small neighborhood, the two given norm germs are
exactly norms defined by weighted filtrations of the reference
residue spaces $V_v=\CC^n$.

Concretely, such a filtration norm is diagonal with real weights
in a complex basis, extended to the valued space using the
reference frame. Two finite filtrations admit a simultaneously
adapted complex basis, obtained by splitting their intersections.
Their canonical norm geodesic therefore interpolates the weights
in that basis. Multiplication of all weights by a positive
number preserves this locus and scales its metric by that number.
This description concerns the finitely many given germs, not a
neighborhood of all norms.

Write their first-order expansions as
\[
 w_{vi}^+(s)=p_{vi}^+s+O(s^2),\qquad
 w_{vi}^-(s)=p_{vi}^-s+O(s^2).
\]
Scale weights by $1/h$ at $s=-h,h$. The endpoint norms converge
to coefficient norms $k_v^-,k_v^+$ with weights $-p_{vi}^-$
and $p_{vi}^+$. Matching and continuity of the midpoint give
\begin{equation}\label{eq:oppositefirstgerms}
 \midpt(k_v^-,k_v^+)=1,
\end{equation}
where $1$ denotes the norm with zero weights in the reference
frame. In a simultaneously adapted complex basis, the paired
weights therefore sum to zero. The two first-derivative lists
agree with multiplicities, and their filtrations are opposite.
In particular, if
\[
 \begin{aligned}
 F_v^+(\le p)&=\Span\{\hbox{positive-side first weights }\le p\},\\
 F_v^-(\ge p)&=\Span\{\hbox{negative-side first weights }\ge p\},
 \end{aligned}
\]
then their intersections give the canonical decomposition
\begin{equation}\label{eq:vertexdecomposition}
 V_v=\bigoplus_p V_{v,p},\qquad
 V_{v,p}=F_v^+(\le p)\cap F_v^-(\ge p).
\end{equation}
Here the first flag is the sum of blocks of index at most $p$,
and the second is the sum of blocks of index at least $p$.

Let $f_0:V_C\to V_A$ be the reduction of the original arrow
in the reference frames. For $s>0$, nonexpansiveness forces
$f_0$ to preserve the increasing flag $F^+$; for $s<0$ it
forces preservation of the decreasing flag $F^-$.
Indeed a nonzero reduced entry must satisfy $p_A\le p_C$
on the positive side and $p_A\ge p_C$ on the negative side.
Hence $f_0$ preserves every summand in
\eqref{eq:vertexdecomposition}. The object
\begin{equation}\label{eq:actualvertexobject}
 E_p=(V_{C,p}\xrightarrow{f_0}V_{A,p})
\end{equation}
is an actual subrepresentation of the reduction. Its support
position is $(x_0,t_0)=(R(u_0)\tan p,\sinh u_0)$.

Fix this $p$ and one side of the vertex. Its finer apartment
frame induces a frame of the $p$-graded quotient, identified
with $E_p$ by \eqref{eq:vertexdecomposition}. A branch of
slope $\alpha\in\{-a,0,a\}$ through $(x_0,t_0)$ has
\[
 w_\alpha'(u)
 =\arctan\frac{x_0+\alpha(\sinh u-\sinh u_0)}{R(u)}.
\]
Its second derivative at the vertex is
\begin{equation}\label{eq:secondvertexweights}
 r_\alpha=(\alpha-\tanh u_0\tan p)\cos^2p.
\end{equation}
This increases strictly with $\alpha$, so its order is $A,B,C$
on both sides of the vertex. The relevant expansions are
$ps+r_\alpha s^2/2+O(s^3)$; the square has the same sign on
both sides.

Inside this $p$ block, a nonzero reduced arrow entry must obey
$r_{\mathrm{target}}\le r_{\mathrm{source}}$, or its norm
inequality would fail for every small nonzero $s$ on that side.
Therefore the spans of blocks with $r\le r_\alpha$ form a
filtration by subrepresentations of $E_p$. Its successive
quotients have the actual arrow reductions of the incident edges.
To see the latter assertion, coordinates with the same $p$ and
$\alpha$ have exactly the same derivative function and, after
normalization at $u_0$, the same weight function. Entries inside
that block retain their reduction. Allowed entries between
distinct second weights are strictly contracting off the vertex;
entries with positive initial valuation remain strictly
contracting on a sufficiently small neighborhood. The same
observation handles distinct first weights. Thus each graded
quotient is the actual semistable edge object, not merely an
object with the same dimension vector. Omitting zero factors
gives increasing phases $1/6,1/2,5/6$.

Doing this on the two sides gives two filtrations of the same
$E_p$. On each open piece the normalized residue arrow is
constant, giving identifications along an edge. At an artificial
subdivision with one continuing branch, each filtration has
just the one factor $E_p$, identifying the adjacent labels.
Coincident branches are treated with their combined multiplicity.
Finally, the first flags are intrinsic one-sided norm germs;
the second flags are their graded limits after subtracting $ps$
and scaling weights by $2/s^2$. The intersection decomposition
and these limits show independence of presentation, up to
isomorphism of the objects, filtrations and edge identifications.
\end{proof}

These are the local residue-category data described in
\cite[\S4.4]{HKS24}. The proposition does not identify them with a
specified global Novikov-field Fukaya object or a Maurer--Cartan
solution. Its input is the norm matching condition; it does not
prove that abstract vertex filtrations imply that condition.

\subsection{Local variational comparison on open pieces}

\begin{lemma}\label{lem:openpiececonverse}
A finite regular harmonic path is locally the exact-endpoint
relaxed minimum away from finitely many times on each compact
interval. First-order matching is not needed for this assertion.
\end{lemma}

\begin{proof}
Its straight branches make the one-sided weight functions
analytic beyond their finite endpoints. Apply the analytic
common-basis assertion of Lemma~\ref{lem:finitecommonbases}
to the two vertex norms on each piece. Refine further at the
zeros of the analytic coordinate gaps $g_i=w_{Ci}-w_{Ai}$.
Each gap on a remaining open piece is either identically zero
or strictly positive. Both vertices now use the same ambient
basis. Put $P_{vi}=R\tan w_{vi}'$.

A zero gap gives a reduced identity arrow on that coordinate.
Its semistable block must have type $B$, so
\[
 w_{Ai}=w_{Ci},\quad P_{Ai}=P_{Ci},\quad
 P_{Ai}'=P_{Ci}'=0.
\]
At a positive gap the reduced source and target are $C$ and
$A$ summands. Semistability of their full velocity blocks
therefore forces
\[
 P_{Ai}'=-aR,\qquad P_{Ci}'=aR.
\]
These statements use the actual coordinate arrows in the
common basis.

On a compact subinterval $[l,r]$ of this piece, first compare
with an ordered curve $z$ in the same apartment and with the
same endpoints. Write $\delta_{vi}=z_{vi}-w_{vi}$.
The Bregman identity and integration by parts give
\begin{equation}\label{eq:scalarcalibration}
 \begin{split}
 E(z)-E(y)={}&\frac12\sum_{v,i}[P_{vi}\delta_{vi}]_l^r
   +\frac12\int_l^r R\sum_{v,i}D_\ell(z_{vi}';w_{vi}')\,\dd u\\
 &+q\sum_{i:g_i\equiv0}\int_l^r R(z_{Ci}-z_{Ai})\,\dd u.
 \end{split}
\end{equation}
There is no basis-volume offset in the gap because both
vertices use the same frame. On positive-gap coordinates the
momentum derivatives cancel the potential; on zero-gap
coordinates they leave the last displayed term.
The endpoint term vanishes, and the other terms are nonnegative.
This proves the comparison for finite action; the infinite
case is immediate. A strict margin is needed only for the
reference curve in the linear term.

For an arbitrary ordered comparison polygon, project its nodes
onto this common apartment at both vertices by the line-sign
projection. It preserves order, endpoints and determinant gap,
and decreases the spectral kinetic action. Apply
\eqref{eq:scalarcalibration} to the resulting polygon.
Passing to all-norm recovery sequences and using regular-curve
recovery proves that $y$ is the local relaxed minimum.
\end{proof}

\subsection{The sign of the boundary terms}
We record the needed CAT(0) fact with its normalization. For an
apartment ray $\gamma(s)$ from $o$ with logarithmic velocity
$d=(d_i)$, put $D=\norm d_2$. Its speed-scaled Busemann function is
\[
 B_\gamma(k)=D\lim_{s\to\infty}
                    \bigl(d_2(k,\gamma(s))-sD\bigr).
\]
For $D=0$ set this function to zero. The limit exists by the
triangle inequality and is finite. If $k$ belongs to the same
apartment, with weight difference $\delta$ from $o$, the
Euclidean distance formula gives
\begin{equation}\label{eq:apartmentbusemann}
 B_\gamma(k)=-\sum_i d_i\delta_i.
\end{equation}

\begin{lemma}\label{lem:oppositebusemann}
Suppose two apartment rays $\gamma_-,\gamma_+$ from $o$ have
opposite weighted germs and the same speed $D$. Then
\begin{equation}\label{eq:busemannsum}
 B_{\gamma_-}(k)+B_{\gamma_+}(k)\ge0
\end{equation}
for every norm $k$. At a matched join of a finite regular path,
the outward momentum germs $-R(u_0)\tan p^-$ and
$R(u_0)\tan p^+$ satisfy this hypothesis, including repeated
velocities and the scalar determinant direction.
\end{lemma}

\begin{proof}
Opposite germs form a local geodesic through $o$. The
concatenation of the two rays is a local geodesic everywhere,
and hence a geodesic line in the CAT(0) norm space.
Thus $d_2(\gamma_-(s),\gamma_+(s))=2sD$.
The triangle inequality through $k$, followed by the defining
limit, proves \eqref{eq:busemannsum}.

For the final assertion, the proof of
Proposition~\ref{prop:vertexfiltrations} identifies opposite
first-derivative flags in one common residue frame.
The odd strictly increasing function $p\mapsto R(u_0)\tan p$
preserves their order, repetitions and opposite weights.
Their momentum rays therefore have opposite germs and equal
Euclidean speed. The finite residue-matrix argument used there
makes this an exact local geodesic assertion for the given
apartment rays. No strict Weyl chamber or trace-free
normalization is required.
\end{proof}

\subsection{The converse across all joins}

\begin{theorem}\label{thm:finiteintervalconverse}
A finite regular harmonic path satisfying first-order matching
is the unique exact-endpoint relaxed minimum on every compact
interval. Competitors may change apartments and may take values
in the complete all-norm space.
\end{theorem}

\begin{proof}
Let $y$ be the path on $[a,b]$. First let $z$ be an ordered split
polygon of finite action with the same full endpoints.
Refine the pieces of $y$ and $z$ at the finitely many exceptions
in Lemma~\ref{lem:openpiececonverse}. At each vertex separately,
apply the analytic common-basis lemma to $y_v$ and $z_v$ on
every closed piece. The weights of $y$ are analytic past the
endpoints, and those of $z$ are affine in its original frame.
A further finite subdivision therefore gives common fixed bases
for $y_v,z_v$, valid at the endpoints as well as in the interior.
The two quiver vertices need not use the same basis.

Write their weights in these bases as $w_{vi},z_{vi}$ and set
\[
 \delta_{vi}=z_{vi}-w_{vi},\qquad P_{vi}=R\tan w_{vi}'.
\]
For $0\le s\le1$, canonical pointwise interpolation $y_s$
between $y$ and $z$ has weights $w_{vi}+s\delta_{vi}$ in
these bases. It is ordered by convexity of the norm order.
It is a finite regular comparison curve, with
\[
 E(y_s)\le(1-s)E(y)+sE(z).
\]
The finitely many compact velocity margins justify
differentiating at $s=0$:
\begin{equation}\label{eq:globalfirstvariation}
 \left.\frac{\dd}{\dd s}\right|_{0+}E(y_s)
 =\int_a^b L(u)\,\dd u,\qquad
 L=\frac12\sum_{v,i}P_{vi}\delta_{vi}'
                   +qR\sum_i(\delta_{Ci}-\delta_{Ai}).
\end{equation}

On an open cell $I$ put $H=\frac12\sum_{v,i}P_{vi}\delta_{vi}$.
Take a nonnegative smooth $\phi$ compactly supported in $I$,
and replace the constant interpolation parameter by
$\varepsilon\phi(u)$. For small $\varepsilon\ge0$ this stays
in $[0,1]$, preserves order, and has a strict velocity margin.
It equals $y$ outside a compact subset of $I$. Local minimality
from Lemma~\ref{lem:openpiececonverse} gives
\[
 0\le\int_I(\phi L+\phi'H)\,\dd u
       =\int_I\phi(L-H')\,\dd u.
\]
The integrand is continuous, so $L-H'\ge0$ on the cell.
This uses an admissible variation of the pair and does not
require the velocity orders at the two vertices to agree.

At an interior subdivision time $u_0$, let
$o_v=y_v(u_0)$ and $k_v=z_v(u_0)$. In the two respective
common endpoint apartments choose rays from $o_v$ with outward
momentum velocities
\[
 -P_v(u_0-),\qquad P_v(u_0+).
\]
They have opposite germs by Lemma~\ref{lem:oppositebusemann},
also at artificial refinements inside an original piece.
Both $o_v$ and $k_v$ belong to each corresponding apartment.
Formula~\eqref{eq:apartmentbusemann} consequently gives the exact
identities
\[
 B_{v,-}(k_v)=\sum_iP_{vi}(u_0-)\delta_{vi}(u_0-),\qquad
 B_{v,+}(k_v)=-\sum_iP_{vi}(u_0+)\delta_{vi}(u_0+).
\]
The left and right bases can differ. The Busemann inequality,
not equality of the coordinate expressions, gives
\begin{equation}\label{eq:joinboundarysign}
 H(u_0-)-H(u_0+)
   =\frac12\sum_v\{B_{v,-}(k_v)+B_{v,+}(k_v)\}\ge0.
\end{equation}

Integrate $L\ge H'$ over every cell and sum. The outer boundary
terms vanish by equality of the full endpoint norms, and
\eqref{eq:joinboundarysign} makes all internal contributions
nonnegative. Thus $\int_a^bL\ge0$. Convexity in $s$ and
\eqref{eq:globalfirstvariation} imply
\[
 E(z)-E(y)\ge
 \left.\frac{\dd}{\dd s}\right|_{0+}E(y_s)\ge0.
\]
This proves comparison with every split polygon.

For a finite-action polygon with arbitrary all-norm nodes,
replace it by the minimum on its fixed partition. Its endpoints
are the split norms of $y$, so
Proposition~\ref{prop:finitepartition} supplies a split polygon
of no larger action with those exact endpoints. The preceding
comparison applies. This replacement need not be close to the
original polygon; only its action and endpoints are used.
Every polygon in an exact-endpoint recovery sequence therefore
has action at least $E(y)$. Regular-curve recovery shows that
$y$ attains the relaxed infimum. Strict midpoint stability
then gives uniqueness in the complete relaxed domain.
\end{proof}

\begin{theorem}\label{thm:wholelineconverse}
Every finite regular harmonic path with first-order matching and
the normalized ends \eqref{eq:normalizedends} equals $c_*$ as
an actual full norm path.
\end{theorem}

\begin{proof}
Let $y$ be such a path. Its finite apartment partition gives
one fixed apartment at each vertex on each distant tail.
A coordinate branch there has
\begin{equation}\label{eq:candidatetailvelocity}
 w'(u)=\arctan(\alpha\tanh u+b\sech u),
 \qquad \alpha\in\{-a,0,a\}.
\end{equation}
At the negative end the coordinate weights have finite limits.
If $\alpha\ne0$, the derivative would tend to the nonzero
constant $\arctan(-\alpha)$, contradicting those limits.
Thus all incoming branches have $\alpha=0$ and type $B$.
The reduced arrow has invertible graded blocks, hence is
itself invertible. For two ordered split norms, invertibility
of the reduced identity arrow means that all relative gaps
vanish. Consequently $y_A=y_C$ on this tail.
Equation~\eqref{eq:candidatetailvelocity} and the full incoming
limit give
\[
 d_2(y_v(u),h_B)=O(e^u),\qquad w_{vi}'(u)=O(e^u).
\]
One common apartment can be used for the equal pair.

On the positive tail, $w_{vi}-\epsilon_vF$ has a finite limit.
Its derivative tends to $\arctan\alpha-\epsilon_v\pi/3$,
so $\alpha=\epsilon_va$. The limiting reference norm belongs
to the same apartment, by closedness after scalar normalization.
Integrating the derivative difference gives
\[
 d_2(y_v(u),r_v(u))=O(e^{-u}),\qquad
 w_{vi}'(u)-\epsilon_v\beta(u)=O(e^{-u}).
\]
These rates show that the nonnegative differential action
$J_{\mathrm{diff}}(y)$ is finite and its tails outside
$[-U,U]$ contribute $O(e^{-U})$.

The exact-reference recovery in the proof of
Theorem~\ref{thm:finitealgebraicminimum} applies verbatim:
keep $y$ on $[-U,U]$, bridge to the reference tails over
intervals of length one, and use the references beyond them.
The resulting $y_U\in\calQ$ converge uniformly to $y$, and
$J(y_U)\to J_{\mathrm{diff}}(y)$. In particular $y\in\calY$
and
\begin{equation}\label{eq:candidaterecovery}
 J_*\le\overline J(y)\le J_{\mathrm{diff}}(y).
\end{equation}

Fix any $z\in\calQ$ and choose $V$ beyond its cutoff and the
tail cutoffs of $y$. Construct $z_V$ by keeping $z$ on
$[-V,V]$, using $y$ outside $[-V-1,V+1]$, and bridging
\[
 y(-V-1)\longrightarrow(h_B,h_B),\qquad
 r(V)\longrightarrow y(V+1)
\]
on the two intervening unit intervals. These are ordered
product geodesics in the respective tail apartments.
The same tail estimates, together with $\beta'=O(e^{-2V})$,
give
\begin{equation}\label{eq:competitortailcost}
 J_{\mathrm{diff}}(z_V)\le J(z)+C_y e^{-V}.
\end{equation}
The constant is independent of $z$ once $V$ exceeds its cutoff.

On $[-V-1,V+1]$ the two curves $y,z_V$ have exactly the same
endpoint norms. Theorem~\ref{thm:finiteintervalconverse},
with regular-curve recovery for $z_V$, gives
$E(y)\le E(z_V)$ on this compact interval. Convert this
comparison to truncated $J$ using
\eqref{eq:renormalization}; the endpoint determinant gaps
agree, so their finite corrections cancel. Outside the
interval the two curves are identical. Adding their finite
common tails yields
\[
 J_{\mathrm{diff}}(y)\le J_{\mathrm{diff}}(z_V)
                      \le J(z)+C_y e^{-V}.
\]
Let $V\to\infty$ and take the infimum over $z\in\calQ$.
Together with \eqref{eq:candidaterecovery}, this proves
$\overline J(y)=J_{\mathrm{diff}}(y)=J_*$.
Uniqueness of the whole-line relaxed minimum gives $y=c_*$.
The argument compares finite renormalized quantities and never
subtracts divergent whole-line unrenormalized actions.
\end{proof}

\begin{proof}[Proof of Theorem~\ref{thm:main}]
Theorem~\ref{thm:finitealgebraicminimum} constructs the finite
regular harmonic path, with its full normalized ends and
quadratic midpoint estimate. Proposition~\ref{prop:finitesupport}
gives its balanced support and vertex bound;
Proposition~\ref{prop:vertexfiltrations} supplies the actual
vertex objects, filtrations and edge identifications.
The tail statements are
Propositions~\ref{prop:outgoingtail} and~\ref{prop:incomingtail}.
Theorem~\ref{thm:wholelineconverse} proves uniqueness in the
stated class. Analytic weights follow from
Proposition~\ref{prop:finitehahnregularity} and the finite
transfer, which preserves the weight functions.

Finally, simultaneous $\GL_n(K)$ pullback preserves order,
relative spectra, determinant gaps, the reference curves and
the relaxed action. It carries $\calQ$ and $\calY$ for one
triple onto those for the transformed triple. Uniqueness
therefore makes the constructed path equivariant. This is
equivariance of full norms with their given framing, without
passing to a quotient by change of basis.
\end{proof}
